\BeleskeID{09_2-s004}
% element: 09_2-s004:e04
Donji ulazni prag traži se u oblasti u kojoj je opterećenje omsko, a pogonski tranzistor zasićen. Leva i desna strana ravnoteže struja diferenciraju se po ulaznom naponu, uz uslov da je nagib karakteristike jednak $-1$:
\[k_{n,L}\left(-2V_{Tn,L}-(V_{DD}-V_O)\right)-k_{n,L}(V_{DD}-V_O)=2k_{n,D}(V_I-V_{Tn,D})\]
\[k_{n,L}\left(-2V_{Tn,L}-2(V_{DD}-V_O)\right)=2k_{n,D}(V_I-V_{Tn,D})\]
Dobijena jednačina se rešava zajedno sa početnom ravnotežom struja. To je postupak sa dve jednačine i dve nepoznate koji se ponavlja pri određivanju naponskih pragova. Prag opterećenja je negativan po prirodi ugrađenog kanala.

Sistem polazi od ravnoteže struja
\[1.\quad k_{n,L}(V_{DD}-V_O)\left(-2V_{Tn,L}-(V_{DD}-V_O)\right)=k_{n,D}(V_I-V_{Tn,D})^2\]
